\section{Introduction}
\label{sec:v56-introduction}

A long billiard bridge is a rare event, but its conditional endpoint law
need not be a quadratic record of the orbit it shadows.  We prove that
on a fixed contact collar it retains nonlinear half-line actions, and
that these actions identify the reflecting contacts.  The forward
statement is relative to an exponentially small endpoint twist.  The
inverse statement concerns the actual boundary, including its curvature,
not an arbitrary formal generating function.  These two statements form
the principal mechanism of the paper.

For clarity we first give the conclusion for a supplied periodic polygon.
A \emph{marked polygon} consists of the lifted collision points $q_i$,
their signed tangent and outward-normal frames $t_i,n_i$, the phase labels
and closing translation, with the specular reflection condition imposed.
Write $L_i=|q_{i+1}-q_i|$, $e_i=(q_{i+1}-q_i)/L_i$ and
$\nu_i=e_i\cdot n_i>0$.  Unknown contacts are represented by
\[
 R_i(x)=q_i+\nu_i^{-1}\{t_i x-n_iF_i(x)\},\qquad
 F_i(0)=F_i'(0)=0,\quad F_i''(0)>0.
\]
The measured coordinate $x=\nu_i t_i\cdot(R_i-q_i)$ is a tangent
projection.  Neither the normal graph height nor the curvature is a mark.
Two fixed positive offsets are observed at every physical phase, on a
common square with positive residuals.  Their limiting densities have
the form
\[
 f_{i,d}(u,v)=Z_{i,d}^{-1}B_i^-(u)B_i^+(v)
              \{d-A_i(u)-C_i(v)\},\qquad d=d_1,d_2,
\]
with anchored actions and amplitudes.  The amplitudes are shared across
the two offsets, while the normalizers need not agree.  The physical
momenta are restored in $A_i=S_i^- -p_i x$ and
$C_i=S_i^+ +p_i x$, where $p_i=e_i\cdot t_i/\nu_i$.

\begin{theorem}[Identification from periodic boundary laws]
\label{thm:v66-main}
Fix a clear nongrazing marked polygon with finitely many distinct physical
contact points in a planar dispersing billiard.  The phase-resolved
limiting laws at two separated positive offsets identify every finite
smooth contact jet without supplied curvatures or amplitudes.  Among
analytic contacts they identify the complete contact germs.  These exact
uniqueness conclusions require no closeness between the candidates.
For actual connected closed analytic obstacle boundaries, the same laws
identify the entire labelled obstacle images when the lattice is fixed
and every obstacle is visited by the polygon.

At each analytic contact tuple the complete future-action map is also a
local biholomorphism in a bounded-holomorphic norm on a sufficiently small
complex disc.  A bounded local analytic prior gives conditional H\"older
stability from real-law error on a smaller disc.  At each separately fixed
smooth order, finite-flight density error gives the geometric error
$C_M(\tau^N+\varepsilon)$ on compact positive classes.
\end{theorem}
\begin{proof}
Theorem~\ref{thm:v64-periodic-relative} supplies the relative limit;
Theorem~\ref{thm:v64-physical-law} and
Lemma~\ref{lem:v65-signed-jacobi} put it in the physical measured
coordinates.  Proposition~\ref{prop:v65-two-offset} separates the
unequal endpoint actions.  The global quadratic inverse of
Theorem~\ref{thm:v66-curvature}, followed by the actual signed recursion,
proves Theorem~\ref{thm:v66-global-rigidity}.  The complete analytic
inverse and conditional real-law estimates are
Theorems~\ref{thm:v65-analytic-inverse} and
\ref{thm:v65-real-stability}.  Proposition~\ref{prop:v66-finite-flight}
proves the fixed-order finite-flight assertion.  All these proofs are
included below.
\end{proof}

The curvature step explains why exact identification is not confined to
one analytic inverse neighborhood.  If $k_i=L_i^{-1}$,
$c_i=F_i''(0)$ and $s_i=(S_i^-)''(0)$, the boundary Schur complement is
\[
 s_i=k_i+c_i-\frac{k_i^2}{k_i+c_{i+1}+s_{i+1}}.
\]
Rearranged as an equation for $c$, it is a strict contraction on the
positive orthant after taking positive parts.  Its constant is bounded
using $k_i$ and the observed $s_i$, not unknown curvature values.
The unique positive fixed point determines the signed orbit factors.
At degree $n$ the actual boundary-visit response is then
$(I+T_n)(I-T_n)^{-1}$, with
$(T_n z)_i=\sigma_i^n z_{i+1}$ and $|\sigma_i|<1$.
This identifies each successive contact jet.  A finite low-degree
inverse together with a norm-convergent high-order tail, rather than
coefficient uniqueness alone, gives the analytic norm theorem.

\subsection{The mechanism and the observation models}

There are three logically different passages.  The cofactor identity and
two-ended trace-class comparison control the relative physical law before
conditioning.  The stationary boundary envelope, including its vanishing
terminal term, converts its nonlinear actions into contact geometry.
Registration and analytic continuation then assemble geometry in the
appropriate observation model.  The latter passage does not replace the
first two.  The realized equal-leading-data family below shows why a
quadratic approximation would not answer the first question.

The alternating channel uses less supplied first-order geometry: its gap
and signed transverse conventions replace a full marked polygon, and one
same-type law at a fixed positive offset suffices at each end.  Its
normal two-site response is the specialization of the signed cyclic
response, but its law-to-action extraction is a different theorem.  Its
unregistered multichannel reconstruction can recover an unknown lattice.
In the marked-polygon conclusion above, in contrast, lattice and placement
are supplied, and full visitation turns exact contact identification into
whole-boundary identity.  No implication that removes these supplied
marks is used.  The charged finite-preparation results require their
calibration and sensor bounds in addition to the geometric inverse;
they are not general-period sampling rates inferred from density error.

The nearby spectral questions concern different observations.
B\'alint, De Simoi, Kaloshin and Leguil~\cite{BDKL}, Theorem D and
Corollary E, extract general-period Lyapunov information from marked
length asymptotics; their Remark 2.3 explains the contact-separation
issue for those asymptotics.  Here phase-resolved functions and a supplied
polygon provide individual action responses.  We do not infer them from
a marked length spectrum.  The enriched spectrum of
Finamore--Leguil~\cite{FinamoreLeguil} and the analytic symmetric
open-billiard setting of De Simoi--Kaloshin--Leguil~\cite{DSKL} likewise
have their own data and hypotheses.  The exact global uniqueness proved
here is for the stated marked boundary-law observation.

The periodic relative and local coordinate statements are recorded next.
Their exact local uniqueness clauses are strengthened by
Theorem~\ref{thm:v66-global-rigidity}; their analytic norm estimates retain
their stated neighborhoods, priors and radius loss.  The subsequent
alternating formulations preserve the weaker-data results and their
complete proofs.


\subsection{Relative transmission along an arbitrary periodic itinerary}

The forward mechanism extends to every selected nongrazing periodic
orbit with clear flights and strictly dispersing contacts.  Its proof
uses the same two-ended normalization, now with periodic Jacobi
coefficients and distinct future and past boundary actions.  We give
the geometric verification rather than assuming a hyperbolic normal
form.  A three-obstacle orbit and its nonsymmetric perturbations provide
realizations outside the normal alternating setting
(Proposition~\ref{prop:v64-triangle}).

Here is the normalization which makes the extension precise.  Let $P$
be an even oriented period, $L_i$ the flight lengths, $\kappa_i$ the
curvatures, and $\nu_i$ the outgoing normal cosines.  Set
$k_i=L_i^{-1}$, $m_i=2\kappa_i/\nu_i$ and
\[
 Q_i=\begin{pmatrix}1+m_i/k_i&1/k_i\\m_i&1\end{pmatrix},\qquad
 \mathcal M_b=Q_{b+P-1}\cdots Q_b,\qquad
 \operatorname{tr}\mathcal M_b=2\cosh\chi.
\]
Positive curvature gives $\chi>0$.  For $N=nP$ flights starting at
phase $b$, the exact reference twist is
\[
 D_{N,b}=\frac{\sinh\chi}{(\mathcal M_b)_{12}\sinh(n\chi)}.
\]
In scaled signed arclength coordinates, let $p_b$ be the reference
endpoint momentum and $L=\sum L_i$.  Theorem~\ref{thm:v64-periodic-relative}
proves, on a collar independent of $n$,
\begin{equation}\label{eq:v64-intro-relative}
 \begin{split}
 W_{N,b}-nL+p_bu-p_bv
     &=S_b^-(u)+S_b^+(v)+O_{C^k}(\tau^N),\\
 -W_{N,b,uv}/D_{N,b}
     &=B_b^-(u)B_b^+(v)+O_{C^k}(\tau^N).
 \end{split}
\end{equation}
Each differentiation order is fixed separately.  The future and past
actions need not agree.  The physical window law contains
$d-S_b^-(u)-S_b^+(v)+p_bu-p_bv$; the endpoint momentum terms
cannot be discarded as though the reference collisions were normal.
Theorem~\ref{thm:v64-physical-law} derives this law from normalized
phase volume and retains its exact exponentially small acceptance mass.

The same estimate also identifies the nonlinear coefficient of the weak
interaction between the ends.  If
$C_N=W_N(u,v)-W_N(u,0)-W_N(0,v)+W_N(0,0)$ and
$\zeta_b^\pm(x)=\int_0^x B_b^\pm(t)\,\dd t$, then
\[
 -C_N/D_{N,b}=\zeta_b^-(u)\zeta_b^+(v)+O_{C^k}(\tau^N).
\]
The error before normalization is
$O(D_{N,b}\tau^N|uv|)$.  Moreover $\zeta_b^-$ linearizes the
future stable period map, while $\zeta_b^+$ linearizes the reversed
unstable period map, both with multiplier $e^{-\chi}$
(Corollary~\ref{cor:v64-transmission}).  Scalar linearization and
the final double integration are consequences of the relative theorem;
they are not substitutes for its geometric determinant estimate.

This extension explains a structural feature of the mechanism: it
requires dispersing coercivity, a clear nongrazing itinerary and summable
endpoint layers, rather than equal flight lengths or a normal two-orbit.
The inversion below exploits the additional geometry of an alternating
pair.  All of its complete-jet, analytic-germ, global and finite-observation
conclusions are retained.  No two-contact inverse is silently applied to
the distinct actions of a general periodic itinerary.

\subsection{The periodic law as a geometric coordinate}

The two-ended law has a further inverse consequence beyond normal
alternating channels.  Fix a clear nongrazing periodic polygon with $r$
distinct marked contacts.  Its collision points, tangent lines, phase
labels and closing translation are given; its contact curvatures and
normal graph germs are unknown.  At every physical phase retain the
conditional law of the two scaled tangent projections at two fixed,
distinct positive offsets.  These are function-valued data on a common
interior square, not samples of the unknown normal graph.

\begin{theorem}[Periodic contact reconstruction]\label{thm:v65-intro-inverse}
For the fixed marked polygon, the phase-resolved two-offset laws locally
determine every finite contact jet in the smooth category, including
curvature.  At an analytic base tuple they determine the complete contact
germs, and the future action map is a local biholomorphism in a
bounded-holomorphic norm on a small fixed disc.  Under a bounded local
analytic prior, real uniform density errors give conditional H\"older
bounds on a smaller disc.  If a fixed-lattice analytic table has every
labelled obstacle visited by the polygon, equality of these laws in the
local inverse class determines all of its obstacle images.
\end{theorem}

The separation of the two actions does not assume that they agree.
Writing the limiting density as
\[
 f_d(u,v)=Z_d^{-1}B^-(u)B^+(v)\{d-A(u)-C(v)\},
 \qquad A(0)=C(0)=0,
\]
the quotient $Q=f_{d_1}(u,v)f_{d_2}(0,0)/
[f_{d_2}(u,v)f_{d_1}(0,0)]$ gives
\[
 A(u)+C(v)=\frac{d_1d_2(1-Q)}{d_2-d_1Q}.
\]
Both flux factors and both normalizers cancel.  The oblique linear
momenta are then restored when passing from these physical actions to
the anchored variational actions.

There is also a geometric reason why all contact orders are invertible.
Let $\sigma_i$ be the signed one-step derivative of the future stable
half-line in normalized graph coordinates.  Dispersing positivity gives
$|\sigma_i|<1$.  At contact order $n$, the response is
\begin{equation}\label{eq:v65-intro-cayley}
 \mathcal C_n=(I+T_n)(I-T_n)^{-1},\qquad
 (T_nz)_i=\sigma_i^n z_{i+1}.
\end{equation}
This formula is obtained by counting the initial boundary visit once and
each internal visit twice in the actual stationary envelope.  It is not
an identification of a formal series with a smooth action.  The block
has determinant one for an even cycle; the explicitly computed odd-cycle
determinant is different but nonzero, including at odd contact orders.
The derivative of the full action is an invertible multiplier plus
contracted boundary evaluations.  A finite low-degree inverse and one
operator-norm tail inverse control all its lower-degree couplings.

\begin{proof}[Proof of Theorem~\ref{thm:v65-intro-inverse}]
Proposition~\ref{prop:v65-two-offset} separates the physical actions.
Lemma~\ref{lem:v65-envelope} proves smooth-jet factorization, and
Theorem~\ref{thm:v65-cyclic-jets} proves \eqref{eq:v65-intro-cayley}
and the local finite-jet inverse, beginning with its invertible
curvature block.  Lemma~\ref{lem:v65-holomorphic} constructs the
actual Banach-valued action map.  Its complete inverse, real-observation
estimate and exact analytic continuation consequence are respectively
Theorem~\ref{thm:v65-analytic-inverse},
Theorem~\ref{thm:v65-real-stability} and
Corollary~\ref{cor:v65-law-rigidity}.  All proofs are given in
Section~\ref{sec:v65-periodic-inverse}.  A nonnormal three-contact
family with independently varying curvatures is constructed in
Proposition~\ref{prop:v65-nonnormal-rigidity}.
\end{proof}

The relationship with variational and spectral methods is specific.
Hill's formula~\cite{Hill} concerns periodic Hessians and monodromy;
our forward normalization uses a Dirichlet corner cofactor before
conditioning.  The general-word length asymptotics and marked Lyapunov
recovery of B\'alint, De Simoi, Kaloshin and Leguil~\cite{BDKL}
concern periodic-orbit lengths, rather than endpoint-varying probability
laws.  Here the fixed marked polygon and the phase-resolved function data
make a separate contact inverse possible.  We do not identify these data
with an ordinary or enriched marked length spectrum.  In particular,
the enriched-spectrum rigidity theorem of Finamore and
Leguil~\cite{FinamoreLeguil} is a different observation model, not a
claimed consequence of the present theorem.  The matrix identity and
the functional-analytic inversion argument are not advanced as separate
novel rigidity theories; their geometric realization in the normalized
physical law and the actual contact action is the point of the argument.

The following alternating-channel results retain a different advantage:
they require only the gap and one law of each contact type, rather than
a fixed general polygon and two offsets per phase.  They also retain
the unknown-lattice reconstruction, finite global ambiguity and charged
finite-preparation results under their original hypotheses.  The new
periodic contact theorem extends the inverse domain without changing
those data requirements or replacing their conclusions.


\subsection{The alternating channel and its geometric content}

A long alternating billiard bridge has two different geometric scales.
Its middle approaches a hyperbolic two-periodic orbit, whereas its
transverse endpoints remain in a fixed contact collar.  Does conditioning
on this increasingly rare event retain the nonlinear geometry of the
reflecting boundaries, or only their contact Hessians?  We prove that it
retains the two signed half-line actions and that these actions determine
every finite contact jet.  For analytic periodic tables, a finite selected
set of marked channels then determines complete obstacle images and the
unknown marked lattice, with finite ambiguity in the noncircular class
and uniqueness in the properly asymmetric class.

The offset and the flight number play different roles.  We let the number
of flights tend to infinity at a \emph{fixed positive excess time}; we do
not replace that positive-offset law by its leading small-offset term.
The unconditioned mass decays exponentially.  The normalized law,
however, retains a nonlinear endpoint interaction on a collar independent
of flight number.  Thus a long-flight limit is not a quadratic truncation.
The separation is realized by actual periodic tables: an analytic family
has the same gap, free area, contact curvatures, and leading endpoint
metrics and count amplitudes at every flight number, but different
nonlinear limiting laws.  Its construction and its nonzero response are
proved in Theorem~\ref{thm:v4-jet-fiber}.

There are two analytical steps.  First, one must control the physical
mixed derivative relative to an exponentially small reference twist.
An absolute estimate for the stationary action cannot do this.  An exact
cofactor identity followed by a two-ended trace-class comparison produces
the relative limit before conditioning and before offset differentiation.
Second, one must show that the action jet is a function of the actual
boundary jet, independent of arbitrary smooth remainders.  A finite
stationary-envelope identity, with its terminal term retained until it
vanishes, proves this factorization.  The resulting signed inverse has a
determinant-one block at every order, without reflection symmetry or
equal contact curvatures.  These are the two uses of the boundary-layer
construction on which the reconstruction rests.



\subsection{The local relative and smooth inverse mechanism}

Fix a clear closest channel between two smooth strictly convex reflecting
arcs, with gap $g>0$, positive curvatures $\kappa_0,\kappa_1$, and a
common signed transverse coordinate.  In contact frames the arcs are
$(-\psi_0(u),u)$ and $(g+\psi_1(v),v)$, with
$\psi_b(0)=\psi_b'(0)=0$ and $\psi_b''(0)=\kappa_b$.
The stationary length of $j$ alternating flights is $W_j(u,v)$.
Put $c_b=1+g\kappa_b$, $c=\sqrt{c_0c_1}$,
$\gamma=\operatorname{arcosh}c$, and $p=j\bmod2$.  Geometric
parameter derivatives, when taken, keep the excess time
$d=t-jg$ fixed in these moving intrinsic coordinates.

\begin{theorem}[Relative boundary law and actual smooth contact inverse]
\label{thm:v56-local-mechanism}
For a selected clear channel in a smooth periodic dispersing table,
there are a flight-independent contact collar, smooth strictly convex
actions $S_b$ and positive smooth amplitudes $B_b$, normalized by
$S_b(0)=S_b'(0)=0$ and $B_b(0)=1$, with the following properties.
\begin{enumerate}
\item Write $a_b=S_b''(0)=c\sinh\gamma/(gc_{1-b})$ and
$d_j^0=\sqrt{a_0a_p}/\sinh(j\gamma)$.  At every separately fixed
finite differentiation order, on a sufficiently small common collar,
\begin{equation}\label{eq:v56-local-relative}
 \left\|W_j-jg-(S_0\oplus S_p)\right\|_{C^k}
 +\left\|\frac{-W_{j,uv}}{d_j^0}-B_0\otimes B_p\right\|_{C^k}
                         \le C_k\tau^j,\qquad 0<\tau<1.
\end{equation}
The estimates include the stipulated fixed-order family derivatives;
the constants may depend on higher smooth norms.  No estimate uniform
in the differentiation order is asserted.
\item At fixed sufficiently small positive excess $d$, the limiting same-type
conditional endpoint law obtained from the full phase ensemble has
positive interior density
\begin{equation}\label{eq:v56-local-density}
 f_b(u,v)=Z_b^{-1}B_b(u)B_b(v)\{d-S_b(u)-S_b(v)\}.
\end{equation}
The gap and one such law of each type determine the two signed action
germs and every finite labelled contact jet, without supplied amplitudes,
curvatures or normal collision coordinates.  The extraction is smooth at
each fixed finite order on the positive, nondegenerate domains specified
in the density inverse below.
\item For $n\ge3$, after lower jets are fixed, write
$s_n=(S_0^{(n)}(0),S_1^{(n)}(0))^t$ and
$q_n=(\psi_0^{(n)}(0),\psi_1^{(n)}(0))^t$.  Then
\begin{equation}\label{eq:v56-local-block}
 s_n=M_nq_n+R_n,\qquad
 M_n=\begin{pmatrix}
 \coth(n\gamma)&\mathfrak r_0^n\operatorname{csch}(n\gamma)\\
 \mathfrak r_1^n\operatorname{csch}(n\gamma)&\coth(n\gamma)
 \end{pmatrix},\quad
 \mathfrak r_b=\sqrt{c_b/c_{1-b}},\quad \det M_n=1.
\end{equation}
Here $R_n$ depends only on the gap and lower graph jets.  The triangular
inverse holds for actual smooth representatives, including odd contact
orders.  More precisely, let $M\ge2$ and let two anchored graph pairs
have the same gap and equal graph jets through order $M$ at their
respective anchored contacts.  Under common functional smooth bounds
on a sufficiently small collar, their actions satisfy
\begin{equation}\label{eq:v56-local-remainder}
 |S_b^{[1]}(u)-S_b^{[0]}(u)|
       \le \frac{C_M}{1-\rho^{M+1}}|u|^{M+1},\qquad 0<\rho<1.
\end{equation}
The highest-degree blocks have the additional bounds
\begin{equation}\label{eq:v58-headline-block-bounds}
 \begin{split}
 \max\{\|M_n\|_\infty,\|M_n^{-1}\|_\infty\}
       &\le \frac{3+t^6}{1-t^6},\\
 \max\{\|M_n-I\|_\infty,\|M_n^{-1}-I\|_\infty\}
       &\le \frac{4\vartheta^n}{1-t^6},
 \end{split}
 \qquad t=e^{-\gamma},\quad \vartheta=\frac{1+t^2}{2}<1.
\end{equation}
Here the norm is the maximum absolute row sum.  These bounds are uniform
in $n\ge3$, and uniformly over geometries with $\gamma\ge\gamma_*>0$.
They do not assert an order-uniform bound for the lower-jet remainders
or the complete nonlinear inverse.
\end{enumerate}
The local forward and finite-jet conclusions do not require analytic
boundaries, equal curvatures, or reflection symmetry.
\end{theorem}

\begin{proof}
Lemma~\ref{lem:g-relative} constructs the finite stationary bridges.
Lemma~\ref{lem:v4-halfline} constructs their half-line limits.
The two-ended normalized determinant comparison in
Theorem~\ref{thm:v4-factorization} proves
\eqref{eq:v56-local-relative}.  The full phase formula
\eqref{eq:g-full-flux}, the common-domain integration in
Lemma~\ref{lem:g-radial}, and Theorem~\ref{thm:v4-law} give
\eqref{eq:v56-local-density}.  Theorem~\ref{thm:v26-density-inverse}
and Proposition~\ref{prop:v26-density-stability} extract the actions and
leading curvatures on the positive interior domain.
Lemma~\ref{lem:v27-smooth-jet-factorization} proves
\eqref{eq:v56-local-remainder} before the last-jet calculation is made;
Proposition~\ref{prop:v22-last-jet-block} gives
\eqref{eq:v56-local-block} and its inverse.
Proposition~\ref{prop:v58-uniform-blocks} uses the positive-curvature
constraint to prove \eqref{eq:v58-headline-block-bounds}.
All of these supporting proofs are printed in the present article.
The block bounds are order-uniform; the nonlinear smooth estimates
and reconstruction remain at each separately prescribed finite order.
\end{proof}

\subsection{Why the relative estimate is the relevant limit}

The estimates have different roles.  To prove the relative estimate one
uses, before division by $d_j^0$, the exact identity
\begin{equation}\label{eq:v56-normalized-determinant}
 \log b_j=\sum_{i=0}^{j-1}
 \log\{g[-\ell_{i\bmod2,uv}(y_i,y_{i+1})]\}
                         -\log\det(I+G_j\Delta H_j),
 \qquad b_j=-W_{j,uv}/d_j^0.
\end{equation}
The perturbation $\Delta H_j$ is localized near two ends; its trace norm
is estimated by a summable weight, not by the number of interior sites.
The finite Green compressions are then compared to their half-line
limits.  Fixed powers introduced by differentiation are absorbed by
strict exponential margins.  For the inverse, one interpolates actual
smooth graphs and differentiates finite stationary action sums.  Interior
variations cancel by stationarity; the terminal term is retained until
its decay has been proved.  Only then is the finite Taylor recursion
used.  Thus neither an absolute-error estimate nor a formal power series
is substituted for the required nonlinear argument.

The distinction can be stated directly at the level of the observation.
On a fixed positive interior square, integrating out the residual time
and normalizing gives
\[
 f_{j,b,d}(u,v)=
 \frac{b_{j,b}(u,v)\{d-E_{j,b}(u,v)\}}
 {\int (d-E_{j,b})_+b_{j,b}}.
\]
The exponentially small reference twist cancels from this quotient.
That cancellation is useful only after $b_{j,b}$, not merely $E_{j,b}$,
has been controlled.  The relative estimate supplies the amplitude and
normalization bounds needed for the positive-offset limit.  In particular,
Corollary~\ref{cor:v26-finite-flight-inverse} gives, at each fixed order
$M$, the geometric error $C_M(\tau^j+\epsilon+\epsilon_g)$ from a
finite-flight interior $C^M$ density error $\epsilon$ and gap error
$\epsilon_g$.  This connects the relative limit to its inverse without
an assumption of order-uniform real differentiation or a bound on the
number of preparations required to estimate the density.

\subsection{The action as a nonlinear boundary coordinate}

The geometric parameter restriction also identifies where an
all-order conditioning issue can occur.  Positivity gives
$\mathfrak r_b e^{-\gamma}<(1+e^{-2\gamma})/2<1$.
Thus the highest-degree response approaches the identity rather than
becoming singular with order.  Proposition~\ref{prop:v58-uniform-blocks}
places this statement on coefficient spaces and for actual smooth
pairs with common lower jets.  The remaining lower-degree couplings,
smooth remainder bounds and analytic continuation are distinct parts
of the full inverse; no estimate for them follows merely from the
diagonal bounds.

The analytic germ problem admits a further conclusion for the complete
nonlinear map, not just its diagonal blocks.  In
Theorem~\ref{thm:v59-analytic-contact-inverse}, the half-line actions
are local analytic coordinates for anchored analytic contact pairs in a
bounded-holomorphic norm on one small fixed disc.  The derivative is
\[
 (D\mathscr S(\psi)\eta)_b(u)
 =\alpha_b(u)\eta_b(u)+
       \sum_{i\ge1}v_{b,i}(u)\eta_{b+i}(x_{b,i}(u)),
       \qquad |x_{b,i}(u)|\le\rho^i|u|,\quad \inf|\alpha_b|>0.
\]
This identity follows from the actual stationary envelope, including its
vanishing terminal term.  The first term is an invertible multiplier;
the remaining evaluations contract high-order vanishing functions.
One finite low-jet system and a norm-convergent tail inverse therefore
control all lower-degree couplings.  A local contraction gives the
nonlinear inverse.  In the induced analytic quotient norms, its complete
linearized finite-jet inverses have bounds independent of the truncation
order.  The analytic norm is an additional hypothesis: it is not inferred
from finitely many density derivatives, total variation, or a finite
sample.  No continuation to an entire obstacle or uniform acquisition
rate is part of this local analytic estimate.

There is also a conditional conclusion in the real observation topology.
Theorem~\ref{thm:v60-real-observation-stability} fixes a local bounded
analytic neighborhood and recovers the complete contact germs on a
strictly smaller complex disc from uniform real endpoint-density error
$\epsilon$, with a bound $C(\epsilon_g+\epsilon^\vartheta)$,
$0<\vartheta<1$.  Curvatures may vary and are not separately supplied.
With a uniform first-order real density bound, total variation or
histogram errors give exponent $\vartheta/3$.  The proof combines
amplitude-free extraction on a real square, a displayed interpolation
estimate, and the complete analytic inverse on the inner disc.
It does not infer same-disc analytic smallness from real data: the
outer analytic prior, radius loss, density and anchor margins, and
local choice of inverse branch are explicit.  No global registration,
unknown-origin calibration, or preparation budget is a consequence of
this local estimate.  Proposition~\ref{prop:v60-contracted-criterion}
isolates the elementary functional-analytic criterion; the protected
half-lines and stationary envelope remain its billiard-specific inputs.


The analytic inverse is another consequence of the stationary envelope,
not a consequence of the determinant of each finite block alone.
The scalar restriction estimate and the interior interpolation used in
the real-observation theorem have a different role: they transfer that
inverse to a weaker observation topology under stated priors.  They are
not additional mechanisms for reconstructing a billiard boundary.
In particular, the all-order conditional estimate and the fixed-order
finite-flight estimate above answer different questions and retain their
different hypotheses.

The density algebra has a separate, simpler role.  On a contact-centered
positive square,
\[
 1-\frac{f(u,v)f(0,0)}{f(u,0)f(0,v)}=t(u)t(v),
 \qquad t(u)=\frac{S(u)}{d-S(u)}.
\]
A nonzero scalar anchor fixes the positive normalization of $t$, and
$S=dt/(1+t)$ recovers the signed action.  This eliminates unknown
separate factors; it is not a new general association invariant.
The signed finite-jet inverse is what turns that invariant into
boundary geometry.  Equality of all smooth jets alone is not identified
with equality of arbitrary smooth germs.  Analyticity enters in the
next step, to recover complete obstacle images.


For fixed contact conventions, the sequence of recoveries is therefore
\[
 (g,f_0,f_1)\longmapsto(g,S_0,S_1)
       \longmapsto (j_0^M\psi_0,j_0^M\psi_1),\qquad M\ge2.
\]
The first arrow is an explicit operation on the interaction of the law;
the second is the inverse of a well-defined map on actual smooth jets.
On an analytic inverse neighborhood the second arrow recovers the entire
contact pair.  This formulation does not identify arbitrary smooth germs
with their jets, nor does it assert that every pair of observed densities
is realized by a billiard.  The theorems reconstruct within the specified
physical class.  For the general marked-polygon, two-offset model,
Theorem~\ref{thm:v68-smooth-rigidity} additionally identifies complete
smooth germs by returning to the full stationary envelope.  That
functional step is distinct from identifying smooth functions with jets;
its finite-smoothness stability is Theorem~\ref{thm:v68-smooth-stability}.

\subsection{The observation and the main theorem}

A table $T$ consists of $N$ labelled strictly convex analytic obstacles
$K_1,\ldots,K_N$ and their translates by $L\mathbb Z^2$, where
$L\in\mathrm{GL}^+(2,\mathbb R)$ is unknown.  Distinct lifted closures
are disjoint and all curvatures are positive.  The basis of the abstract
deck group $\mathbb Z^2$ is marked; its Euclidean realization is not.
A selected clear channel $e=(a,b,\ell)$ joins the closest pair of
$K_a$ and $K_b+L\ell$.  Each channel has its own oriented contact frame
and signed transverse coordinate, with no inter-channel registration.
The full geometric conventions are stated in
Section~\ref{sec:geometric-results} and
Section~\ref{sec:v24-rank-two-lattice}.

For each channel retain its gap $g_e$ and, for each starting contact type
$b\in\{0,1\}$, the limiting conditional joint law $P_{e,b}$ of the two
signed transverse endpoints of an even bridge at one known positive
excess time $d_e$.  The intrinsic datum is
\begin{equation}\label{eq:v48-datum}
       \mathcal D_E(T)=\bigl(g_e,P_{e,0},P_{e,1}\bigr)_{e\in E}.
\end{equation}
Obstacle, channel and deck marks and the signed conventions are design
information.  No longitudinal position, residual time, waiting count,
curvature, flux amplitude, Euclidean deck vector or inter-channel pose
is supplied.  A law is a function-valued datum, not one scalar sample.
The gap is retained separately from the conditional laws.

Let $\mathscr C_N$ be the class in which every obstacle is noncircular,
and let $\mathscr A_N\subset\mathscr C_N$ be the class in which every
obstacle has trivial proper Euclidean symmetry.  The latter is open and
dense in the relative $C^2$ analytic-support topology
(Proposition~\ref{prop:v45-generic-asymmetry}).  A
\emph{common-strip analytic-support family} is $C^1$ into the Banach space
of bounded holomorphic support functions on one fixed complex strip,
with $C^1$ lattice columns; see \eqref{eq:v48-support-family}.
In particular the parameter derivative is itself analytic on a smaller
strip.  The differential statements below do not merely assume that
each parameter slice is analytic.

\theoremstyle{plain}
\newtheorem{structuraltheorem}{Theorem}
\renewcommand{\thestructuraltheorem}{\Alph{structuraltheorem}}
\begin{structuraltheorem}[Single-offset periodic rigidity]
\label{thm:v48-main}
For each $T_0\in\mathscr C_N$ there is a persistent selected design
$E$ of $N+1$ channels, consisting of an obstacle spanning tree and two
independent marked cycle gains, and a common sufficiently small positive
offset, with the following properties.
\begin{enumerate}
\item The datum $\mathcal D_E(T_0)$ determines a finite set of complete
labelled periodic tables, including their marked lattices, modulo one
common proper Euclidean motion.  This set is obtained by finite choices
of incidence congruences followed by explicit lattice reconstruction
and geometric compatibility tests.  If $T_0\in\mathscr A_N$, it
consists of one table.
\item For a common-strip analytic-support variation through $T_0$, if
all gap derivatives and all endpoint-density derivatives on
contact-centered positive interior squares vanish, the table variation
is one common infinitesimal proper Euclidean motion.  Conversely such
a motion has zero intrinsic derivative.  In a fixed Euclidean gauge
the data derivative is injective.
\item On every immersed $p$-dimensional $C^1$ common-strip
analytic-support model in that gauge, $p\ge1$, one can select $p$
scalar functions, each a gap or an expectation
$\int\phi\,dP_{e,b}$ with $\phi\in C_c^\infty$ supported in a positive
interior square, which form local $C^1$ coordinates at $T_0$.
Their inverse is locally Lipschitz, with local control of every fixed
real support norm and of the lattice matrix.
\end{enumerate}
\end{structuraltheorem}

\begin{proof}
Theorem~\ref{thm:v45-clear-skeleton} supplies the persistent design.
The signed law inverse gives complete analytic channel-frame images.
Theorem~\ref{thm:v49-finite-fiber} enumerates the finite matching
alternatives and reconstructs their lattices; the asymmetric uniqueness
conclusion is also Theorem~\ref{thm:v45-generic-determination}.
Theorems~\ref{thm:v48-infinitesimal} and
\ref{thm:v48-finite-coordinates} prove the second and third assertions.
Their local registration uses Lemma~\ref{lem:v49-local-alignment}, not
a choice of one globally preferred symmetry branch.
\end{proof}

The finite bound in the first assertion is a product of orders of
obstacle rotation groups along a spanning tree of channel incidences
(\eqref{eq:v49-fiber-bound}); it is not a sharp degree assertion.
The $N+1$ count is minimal for the stipulated connected two-cycle
architecture, not an information lower bound over arbitrary experiments.
The selected design and its marks are not discovered by an unmarked
trajectory algorithm.  The scalar coordinates in the third assertion
may depend on the base table and model; finite dimensionality is not a
restriction on the first two assertions.

\subsection{Unknown origins and separate recording efficiencies}

The contact-centered formulation isolates the geometric inverse.  Its
origins can themselves be recovered from the law, even when its two endpoint
amplitudes are independently distorted.  Suppose the recorded coordinates
are $x=u+\xi$ and $y=v+\eta$, with unknown constants $\xi,\eta$, and
successful records are retained with a positive factor $e_1(x)e_2(y)$.
The physical signed units and directions are unchanged; neither an unknown
clock nor an unknown coordinate scale is included in these constants.
Section~\ref{sec:v51-interaction} proves the following extension using the
interaction of the two endpoints.

\begin{structuraltheorem}[Rigidity with unknown transverse origins and weights]
\label{thm:v51-main}
In the geometric setting of Theorem~\ref{thm:v48-main}, retain the same
marks, gaps and known positive offsets, but replace every ideal same-type
law by its translated law with unknown positive smooth separate endpoint
recording factors.  Assume that the complete positive cap is recorded and
that the numerical coordinates have the specified signed transverse units
and directions.  Then the laws determine their two transverse origins and
the signed action germs.  The compatible table fiber is exactly the one in
Theorem~\ref{thm:v48-main}(1), including the unknown marked lattices.
In a differentiable table-and-recording family of the stated smooth class,
no table variation outside the common Euclidean tangent space can be
masked by varying those origins or separate recording factors.
\end{structuraltheorem}

\begin{proof}
Lemma~\ref{lem:v51-origins} extracts the origins from the zero fibers of
$\partial_x\partial_y\log\widetilde f$.
Theorem~\ref{thm:v51-action} then recovers the action by its scalar-anchored
four-density ratio.  Proposition~\ref{prop:v51-stability} justifies the
finite-order differentiable extraction, and
Theorem~\ref{thm:v51-geometric} proves the exact-fiber and joint derivative
kernel assertions through the smooth contact inverse and periodic cochain.
\end{proof}

The cancellation underlying this extension is a classical association
invariant~\cite{HollandWang1987,Osius2009}.  What the boundary-law mechanism
supplies is its specific form
\[
 \partial_x\partial_y\log\widetilde f
 =-\frac{S'(x-\xi)S'(y-\eta)}
           {[d-S(x-\xi)-S(y-\eta)]^2}.
\]
Strict convexity locates the two origins, and the signed contact inverse
turns the recovered action into geometry.  No equality of the two recording
profiles is required.  This is an exact-law reconstruction with a local
finite smooth stability estimate, not a claim that a fixed finite vector
of noisy observations or the old preparation budget suffices for arbitrary
unknown recording efficiencies.  Theorem~\ref{thm:v48-main}(3) remains a
separate statement on its stipulated finite-dimensional ideal-law model.





\subsection{The recovered object and the supplied information}

The observation in Theorems~\ref{thm:v48-main} and~\ref{thm:v51-main}
is local in each channel but function-valued.  It supplies the marked
channel design, the signed physical transverse units, gaps and positive
offsets.  It does not supply a lattice Gram matrix, a common pose for
the channel frames, normal collision coordinates, or the separate flux
amplitudes.  No claim that $N+1$ laws constitute $N+1$ numerical samples
is needed for either theorem.  For interior windows, both critical
lines must remain visible; the full cap boundary and discarded mass
are not supplied.

This information permits a precise hierarchy of conclusions.  The
relative and actual-smooth mechanism is local and requires only the
stated smooth bounds.  Analyticity propagates the recovered contact
germs into complete obstacle images.  Noncircularity makes global
congruence choices finite, whereas proper asymmetry makes them unique.
One nonzero support harmonic determines angular velocity along an
actual local alignment branch, so finite symmetry produces no extra
infinitesimal kernel.  The lattice is then recovered from two independent
marked cycle translations using their actual integer gain matrix.
The finite-dimensional scalar-coordinate conclusion is a consequence
of this geometric kernel theorem on a stipulated immersed model.

The nonlinear content is not encoded by the rarity of preparation alone.
The analytic family in Theorem~\ref{thm:v4-jet-fiber} fixes the gap,
free area, contact curvatures and every leading endpoint covariance and
count amplitude, while its fourth contact jet is $3-24s$.  The signed
block inverse and Theorem~\ref{thm:v26-density-inverse} therefore
separate its conditional laws at a fixed sufficiently small positive
offset.  Thus the full law distinguishes geometries that its leading
quadratic record does not.  The relative estimate supplies that nonlinear
record at long flight number; the smooth remainder factorization licenses
its geometric inversion.  These two analytic steps, rather than the
subsequent finite congruence bookkeeping or finite-dimensional coordinate
selection, are the mechanism of the inverse.

For the leading-data comparison the geometric distinction is explicit.
The area-preserving support family in Theorem~\ref{thm:v4-jet-fiber}
has $g=1$, $A=12-\pi$ and $\kappa_0=\kappa_1=1$, while its
fourth graph derivative is $q_s=3-24s$.  The normalized onset response
satisfies
\[
 \left.\frac{d}{ds}\partial_d\mathcal F^{(s)}(0)\right|_{s=0}
                         =\frac{\sqrt3}{2}.
\]
The full proof computes both the action and determinant contributions;
it does not infer this coefficient from an uncontrolled absolute error.
Independently, the signed action inverse separates the fixed-positive-offset
endpoint laws when the fourth contact jets differ.  The example establishes
strict additional geometric information beyond the leading record.  It
does not compare the information of endpoint laws with that of a marked
length spectrum, and its symmetric realization is not an assumption in
the general signed inverse.

The three most direct comparison cases are proved in the article.
Theorem~\ref{thm:v4-jet-fiber} exhibits an analytic family with the same
leading endpoint metrics and count amplitudes but different nonlinear
laws.  Proposition~\ref{prop:v50-two-table} gives a noncircular datum
with exactly two complete periodic realizations and different marked
Gram forms; Proposition~\ref{prop:v50-extra-gap} separates those two
branches with one additional marked gap.  The circular comparison in
Proposition~\ref{prop:v49-circular-comparison} instead exhibits a
continuous lattice ambiguity.  These examples distinguish nonlinear
information, finite exact ambiguity, and an infinitesimal symmetry;
they are not three counterexamples to the structural theorem.

The periodic conclusion addresses a further question after local
reconstruction.  The recovered obstacle images lie in separate channel
frames, not in one supplied Euclidean frame.  Noncircularity makes the
possible incidence congruences finite; the marked cycle gains then turn
the recovered displacement cochain into the unknown lattice.  Thus the
lattice is an output, not hidden metric information in the channel marks.
The congruence and cochain arguments are consequences of the local
reconstruction, rather than independent substitutes for its analytical
content.  Their complete proofs also show why finite rotational symmetry
and a continuous infinitesimal ambiguity must be distinguished.

\subsection{Relation to earlier inverse problems}

Separate-factor cancellation belongs to the association framework of
Holland and Wang~\cite{HollandWang1987} and Osius~\cite{Osius2009}.
The present dynamical input is the limiting density
\eqref{eq:v56-local-density} with its actual nonlinear actions, and the
geometric input is the signed smooth inverse.  These are different
assertions from the algebraic invariance of an odds ratio.

Finamore and Leguil~\cite{FinamoreLeguil} determine finite-horizon Sinai
billiards from an enriched marked length spectrum.  That enriched datum
includes the auxiliary closed-geodesic lengths in their definition and
is not identified here with ordinary periodic-orbit marked lengths.
De Simoi, Kaloshin and Leguil~\cite{DSKL} treat analytic open billiards
under non-eclipse, symmetry and genericity hypotheses.  No reduction
between either marked-length observation and the present endpoint-law
observation is used.  The fifth version of Florio and
Leguil~\cite{FlorioLeguil} removes an earlier geometric spectral-rigidity
assertion affected by an error while retaining dynamical conclusions;
that removed assertion is not a premise here.  Our comparison concerns
what is observed and what is recovered, not an assertion that the present
theorem removes hypotheses from those results.

% Shared hypothesis-accurate comparison, included in both complete entries.
Zelditch's inverse spectral theorem provides a useful orbit-local
comparison~\cite{Zelditch2009}.  Theorem~1.1 of that paper determines a
simply connected analytic plane domain in the class $\mathcal D_{1,L}$
from its Dirichlet or Neumann spectrum.  This class includes an involution
reversing a nondegenerate bouncing-ball orbit, multiplicity one for every
iterated length, and the additional nonresonance and endpoint conditions
specified there.  The proof recovers boundary Taylor coefficients from
localized wave-trace invariants.  The symmetry relates the two boundary
graphs.  Thus neither localization near a periodic orbit nor all-order
boundary recovery, by itself, distinguishes the problem considered here.

The distinction lies in the observation map and in the passage from a
long bridge to its nonlinear endpoint action.  Our datum is a controlled
conditional law obtained from phase-volume preparations, not a Laplace
spectrum or a sequence of wave invariants.  The selected bridge becomes
exponentially rare, but its endpoints remain on a fixed collar at a fixed
positive offset.  Theorem~\ref{thm:v4-factorization} compares the mixed
endpoint derivative with its exponentially small reference value in
\emph{relative} size.  An absolute error bound would not suffice for the
physical conditional normalization.  The resulting law retains the
nonlinear half-line actions; separate-factor cancellation then extracts
those actions from the density.  It is the relative limit, rather than
the cancellation identity alone, that connects an observable law to the
geometric inverse.

The signed contact inverse treats both anchored graphs without imposing
a symmetry between them.  Its finite-order content is a statement about
actual smooth graphs: Lemma~\ref{lem:v27-smooth-jet-factorization} controls
functional remainders by finite-truncation stationarity before passage to
any jet coordinates.  Analyticity enters the subsequent determination of
a complete germ and its continuation, not that smooth factorization.
The function-space inverse controls the full coupled action map, including
its lower-order terms.  With an outer analytic prior, its real-observation
consequence estimates the complete germ on a smaller disc.  The finite
preparation theorems further distinguish the exactly calibrated rate from
the rate that charges the position pilot.  Their minimum-distance
construction is an existence result for a measurable estimator; it does
not supply a finite running-time bound for computational reconstruction.

There is also an internal comparison for which a geometric separation is
proved, rather than inferred from different descriptions of the data.
The family in Theorem~\ref{thm:v4-jet-fiber} has the same gap, area,
quadratic contact data and leading count amplitudes, while its
positive-offset boundary law detects a varying fourth-order contact term.
This shows that the present law retains information absent from those
specified leading data.  It is not a separation theorem between the law
and a full Laplace spectrum, nor between the law and either marked-length
datum discussed above.  No such comparison is needed for the relative
boundary limit or the contact inverse.  These distinctions locate the
mathematical contribution in the nonlinear relative-action mechanism and
its geometric consequences, without treating different observations as
interchangeable.


\subsection{Proof organization and the complete technical manuscript}

Part I contains the full geometric proof: the phase ensemble and finite
bridge, the relative boundary limit, the signed contact inverse,
interior interaction, analytic continuation, and periodic registration.
The differential theorem follows this same chain.  The appendices supply
the graph-to-support conversion, marked source identity and the complete
local comparison inverses invoked in the discussion.  The principal
proof does not use the later stopped-record probability reductions.

The accompanying \emph{complete technical manuscript} retains, in full,
the additional local experiments, the finite-resolution inverse, physical
calibration, and every historical theorem and auxiliary proof.  References
prefixed \textup{F.} point to its independently numbered statements.
They are explicit references to that companion, not missing internal
numbers.  One has a further logical role:
Corollary~\ref{cor:v45-physical-chart} is an application of
Theorem~\ref{thm:v25-global-physical-reconstruction} of the complete
technical manuscript.  It uses that theorem's acquisition and risk
estimates in addition to the geometric results proved here.  Its proof
verifies the persistent-design hypotheses; it does not reprove the
physical pilot or capped estimation theorem.  This application is not
a premise of Theorem~\ref{thm:v56-local-mechanism} or
Theorems~\ref{thm:v48-main} and~\ref{thm:v51-main}.
The companion's quantified acquisition theory uses additional
analytic and sensor bounds; it is not a statistical conclusion under the
hypotheses of Theorem~\ref{thm:v48-main}.  Its successful-mark, waiting-count
and capped-bit experiments are different records of a forward measure,
not additional principal rigidity theorems.  Their complete proofs remain
available without placing those likelihood calculations on the proof
path of Theorem~\ref{thm:v56-local-mechanism} and Theorems~\ref{thm:v48-main} and~\ref{thm:v51-main}.
